\section{Scaling the Harness: Multi-Agent Orchestration over Code}
\label{sec:mas}

As AI systems tackle increasingly complex problems from
function-level code synthesis to repository-level system
engineering, fundamental limitations for single-agent emerge: (1)
context window constraints prevent a single agent from holding an
entire codebase, long interaction history, and execution trace in
working memory; (2) specialization requirements make it
inefficient to use one generalist agent for planning, synthesis,
testing, review, and debugging simultaneously; and (3) the
absence of independent coordination and verification channels
prevents the agent from reliably detecting and correcting its own
errors during long-horizon execution.
Multi-agent systems introduce a powerful principle: once these
responsibilities are distributed across specialized roles, the
agent harness itself becomes more modular, inspectable, and
adaptable.
Early systems such as ChatDev~\cite{Qian2023ChatDev},
MetaGPT~\cite{Hong2023MetaGPT}, and AgentCoder~\cite{huang2023agentcoder}
demonstrate this shift by dividing software-development
responsibilities among distinct agents such as architect,
programmer, tester, reviewer, and executor.
Coordinated through structured communication protocols and shared
code artifacts, these role-specialized agents turn code from a
mere output target into the shared substrate through which the
overall harness plans, acts, verifies, and improves itself.

In this section, we systematically survey the rapidly growing
direction on using MAS to scale coding harnesses, and we propose
a new position on building shared code-centric harness substrates
for AI agents.

\begin{figure}[t]
    \centering
    \includegraphics[width=1.0\linewidth]{figures/mas_harness_overview.pdf}
    \caption{Overview of scaling the agent harness through multi-agent orchestration over code. The figure illustrates how role-specialized agents, shared code-centric substrates, execution feedback, and adaptive collaboration topologies address single-agent limitations in context, specialization, and self-correction.}
    \label{fig:mas-overview}
\end{figure}

\begin{figure}[t!]
    \centering
    \includegraphics[width=0.85\linewidth]{figures/mas_roadmap_new.pdf}
    \caption{Roadmap of scaling code harnesses for multi-agent orchestration, organized by workflow collaboration, shared repository state, execution verification, and adaptive coordination.}
    \label{fig:roadmap_sec4}
\end{figure}

\subsection{Improved Coding Support through Multi-agent Collaboration}

The most immediate contribution of multi-agent systems is that
they improve coding support by decomposing the harness into
specialized but coordinated components. Instead of integrating
planning, synthesis, execution, and verification into a single
agent loop, these systems distribute responsibility across roles
that interact through shared code artifacts and feedback signals.
This division of labor makes the overall harness more capable of
handling complex software tasks, while also making its internal
workflow more inspectable and controllable. In practice, this
improvement is realized through three closely related design
dimensions: how roles are specialized, how agents interact over
shared program artifacts, and how the workflow topology organizes
their collaboration.




\begin{table}[t]
\centering
\caption{Representative MAS collaboration designs by role specialization and interaction structure. }
\label{tab:mas_collaboration_systems}
\renewcommand{\arraystretch}{1.16}
\setlength{\tabcolsep}{3.5pt}
\footnotesize
\begin{tabularx}{\textwidth}{@{}
  >{\raggedright\arraybackslash}p{2.55cm}
  >{\raggedright\arraybackslash}p{2.85cm}
  >{\raggedright\arraybackslash}p{3.05cm}
  >{\raggedright\arraybackslash}p{3.05cm}
  >{\raggedright\arraybackslash}X@{}}
\toprule
\textbf{System}
& \textbf{Harness Substrate}
& \textbf{Agent Roles}
& \textbf{Interaction Mode}
& \textbf{Topology} \\
\midrule
Self-Collaboration~\cite{Dong2024SelfCollaboration}
& Blackboard, implicit
& Plan, Synth., Verif. (simulated)
& Critique-repair
& Pre-defined cyclic \\ \midrule
CodePori~\cite{Rasheed2024Codepori}
& Implicit
& Plan, Synth., Verif.
& Collab-Synth., critique-repair
& Pre-defined chain, cyclic \\ \midrule
MAGIS~\cite{Tao2024Magis}
& Repository, evolution memory
& Plan, Understand, Synth., Verif.
& Critique-repair, debate, delegation
& Hierarchical, cyclic, dynamic pool \\ \midrule
HyperAgent~\cite{Phan2024HyperAgent}
& Repository, execution
& Plan, Understand, Synth., Exec
& Critique-repair
& Pre-defined hierarchical, cyclic \\ \midrule
PairCoder~\cite{Zhang2024PairProgramming}
& Execution
& Plan-Understand, Synth-Exec
& Collab-Synth., critique-repair
& Pre-defined cyclic with conditional branch \\ \midrule
FlowGen~\cite{Lin2025Soen101}
& Execution, implicit
& Plan, Understand, Synth., Verif.
& Critique-repair, debate
& Pre-defined chain, cyclic (Scrum) \\ \midrule
Trae Agent~\cite{gao2025traeagent}
& Repository, execution
& Generate, Prune, Select
& Collab-Synth., search (selection)
& Pre-defined search pipeline \\ \midrule
BOAD~\cite{xu2025boad}
& Repository, execution
& Orchestrate, Localize, Edit, Validate
& Delegation, adaptive selection
& Adaptive hierarchical \\ \midrule
FlowReasoner~\cite{gao2025flowreasoner}
& Execution, implicit
& Meta-design, Solve
& Runtime workflow generation
& Objective-driven adaptive \\ \midrule
ChatDev~\cite{Qian2023ChatDev}
& Implicit, borderline exec
& Plan, Synth., Verif., Exec
& Critique-repair, debate
& Pre-defined chain (waterfall) \\ \midrule
MetaGPT~\cite{Hong2023MetaGPT}
& Implicit, partial blackboard
& Plan$\times$3, Synth., Verif.
& Critique-repair, pub-sub scheduling
& Pre-defined chain (waterfall) \\ \midrule
GameGPT~\cite{chen2023gamegpt}
& Blackboard (dual collaboration)
& Plan, Synth., Verif.
& Critique-repair, collaborative
& Pre-defined \\
\bottomrule
\end{tabularx}
\end{table}




\subsubsection{Functional Role Specialization and Human-Guided Planning}
In human software development, different roles specialize in
different aspects of the development process. Many MAS naturally
mirror this division of labor by assigning distinct functional
roles to different agents. This specialization allows each agent
to focus on a specific slice of the shared code harness,
leveraging its unique capabilities and perspectives to contribute
to the overall task.
Here, we elaborate on the most common functional roles identified
across the surveyed literature, noting that many systems
implement multiple roles and that the boundaries between them can
be fluid.

\paragraph{Program synthesis agents} Program synthesis agents are
responsible for generating or transforming code. They consume
specifications, plans, or feedback signals and produce or revise
code artifacts. This is the most common role across surveyed
systems. Representative instances include the Coder in
Self-Collaboration~\cite{Dong2024SelfCollaboration}, the
Programmer in AgentCoder~\cite{huang2023agentcoder}, the Engineer
in MetaGPT~\cite{Hong2023MetaGPT}, the Developer in
ChatDev~\cite{Qian2023ChatDev}, and the RTL Generation Agent in
MAGE~\cite{Zhao2024MAGE}.

\paragraph{Program understanding agents} Program understanding
agents analyze existing code or specifications to produce
higher-level representations. They own the interpretation of what
the code means rather than what it does. This category includes
the Repository Custodian in MAGIS~\cite{Tao2024Magis}, the
Navigator in HyperAgent~\cite{Phan2024HyperAgent}, the RepoUer in
Lingma SWE-GPT~\cite{Ma2024Lingma}, and the Column-type Annotator
in CleanAgent~\cite{Qi2024CleanAgent}.

\paragraph{Verification agents} Verification agents evaluate code
quality, typically by generating test cases, running static
analysis, or simulating execution. The Test Designer in
AgentCoder~\cite{huang2023agentcoder} generates test cases
independently of the code to avoid circular reasoning, a design
principle against the mode-collapse problem where an agent's
biased tests pass its own buggy code. The Test Quality Checker in
QualityFlow~\cite{Hu2025QualityFlow} addresses this at a
meta-level, filtering synthesized tests before they are used as
feedback. The Static Analyzer and Fuzzing Agent in
AutoSafeCoder~\cite{Nunez2024AutoSafeCoder} provide
security-oriented verification through static CWE analysis and
dynamic crash detection, respectively. The Panelists in
CANDOR~\cite{Xu2025Hallucination} independently audit oracle
correctness against natural language specifications rather than
against the code itself, deliberately avoiding contamination by
faulty implementations.

\paragraph{Execution agents} Execution agents interface directly
with the program runtime. Critically, the Test Executor in
AgentCoder~\cite{huang2023agentcoder} is a deterministic Python
script (not an LLM) which cleanly separates reasoning from
execution and grounds the feedback signal in objective program
behavior. The Executor in HyperAgent~\cite{Phan2024HyperAgent}
runs unit and integration tests via an interactive bash shell.
The Judge Agent in MAGE~\cite{Zhao2024MAGE} interfaces with RTL
simulation tools to produce per-clock-edge waveform snapshots.

\paragraph{Planning agents} Planning agents decompose the overall
software-development task into subtasks and assign them to
synthesis agents. The Architect and Project Manager in
MetaGPT~\cite{Hong2023MetaGPT}, the Manager in
MAGIS~\cite{Tao2024Magis}, the Scrum Master in
FlowGen~\cite{Lin2025Soen101}, and the Mother agents in
SoA~\cite{Ishibashi2024SelfOrganized} all perform task
decomposition. The Mother agents in
SoA~\cite{Ishibashi2024SelfOrganized} are particularly notable:
they dynamically spawn Child agents at runtime based on the
inferred complexity of each subfunction, making planning and
agent initialization interdependent.

A distinctive feature of EvoMAC~\cite{Hu2025EvoMAC} is the
introduction of two novel meta-roles not present in any other
surveyed system: the Gradient Agent, which reads execution logs
to identify which agents caused failures, and the Updating Agent,
which revises agent prompts and restructures the workflow DAG
accordingly. These roles operate at the level of the MAS itself
rather than the program, enabling the system to adapt its own
structure in response to execution feedback.


\subsubsection{Diverse Interaction Modes Grounded in Shared Program State}
Unlike general MAS where agent interaction is primarily
message-passing, code-centric interaction is characterized by
artifact-mediated communication: agents observe and modify shared
code artifacts, and their interaction is grounded in the
objective state exposed by those artifacts and their execution
results. {These coordination channels are broader
than source code alone: agents communicate through APIs, files,
diffs, tests, logs, schemas, blackboards, and explicit workflow
states. In most systems, these channels are part of the
human-designed harness, while agents dynamically write to or
modify the artifacts circulating within them.} We identify four
interaction modes.

\paragraph{Collaborative synthesis} Collaborative synthesis occurs
when two agents jointly construct a program component, analogous
to pair programming \cite{zou2026recursivemas}. The Navigator--Driver pairing in
PairCoder~\cite{Zhang2024PairProgramming} is the most direct
instantiation: the Navigator generates and selects solution plans
while the Driver implements them, with bidirectional information
flow. CodePori~\cite{Rasheed2024Codepori} implements
collaborative synthesis between Dev\_01 and Dev\_02, who exchange
code drafts across three rounds. This mode is relatively rare among the surveyed system, as most systems prefer a sequential handoff
rather than true co-construction.

\paragraph{Critique and repair} Critique and repair is the
dominant interaction mode across the surveyed literature. A
verification or evaluation agent inspects a code artifact and
produces structured feedback; a synthesis agent then revises the
artifact in response. This pattern appears in some form in
virtually every surveyed system. Its key design decisions are:
(a) whether the critique is LLM-simulated or execution-grounded
(Self-Collaboration~\cite{Dong2024SelfCollaboration} uses a
simulated LLM tester, while
AgentCoder~\cite{huang2023agentcoder} uses a real Python
executor); (b) the richness of the feedback signal (ranging from
binary pass/fail in SEW~\cite{Liu2025SEW} to structured execution
logs enumerating satisfied requirements, function errors, and
unmet requirements in EvoMAC~\cite{Hu2025EvoMAC}); and (c) the
number of repair iterations permitted before fallback.

\paragraph{Adversarial validation} Adversarial validation is a
more active form of verification in which one agent attempts to
break the code through adversarial inputs, rather than passively
reviewing it. AutoSafeCoder~\cite{Nunez2024AutoSafeCoder}
implements this via its Fuzzing Agent, which generates
crash-inducing input seeds using type-aware mutation and executes
the code to produce crash traces. This mode has a fundamentally
different character from critique-and-repair: the fuzzer does not
explain what is wrong, but demonstrates a concrete execution
failure, a counterexample that the coding agent must address.
MAGE~\cite{Zhao2024MAGE} similarly uses simulation mismatch as an
adversarial signal: the Debug Agent receives the exact waveform
window around the first clock-edge failure, enabling targeted
repair.

\paragraph{Reasoning debate} Reasoning debate involves agents
arguing over the correctness of a decision or the interpretation
of a specification, before arriving at a consensus.
ChatDev~\cite{Qian2023ChatDev} introduces communicative
de-hallucination, a mechanism in which the assistant agent
reverses roles to ask clarifying questions before committing to a
response. The Scrum sprint meetings in
FlowGen~\cite{Lin2025Soen101} enable disordered multi-agent
discussion around a shared context buffer before the Scrum Master
synthesizes a decision. CANDOR~\cite{Xu2025Hallucination}
implements the most structured debate mechanism: three
independent Panelists evaluate oracle correctness, and a Curator
aggregates their verdicts via majority vote. The kick-off meeting
in MAGIS~\cite{Tao2024Magis} involves a circular speech among the
Manager and all Developer agents to negotiate task dependencies
and prevent conflicts.


\subsubsection{Optimized Workflow Topology for Agentic Coordination}

The topology of agent interaction, who communicates with whom, in
what order, and how many times, is one of the most consequential
design decisions in a MAS for code generation. We organize
topologies along two primary axes.

\paragraph{Pre-defined Heuristic Topologies}
The majority of surveyed systems use topologies that mirror
established software development life cycle (SDLC) models. These
topologies are fixed at design time and do not change in response
to task complexity or system performance.

\textbf{\textit{Chain (Waterfall) topologies}} sequence agents in
a strict linear order, with artifacts flowing unidirectionally
from planning to synthesis to verification. ChatDev~\cite{Qian2023ChatDev}
and MetaGPT~\cite{Hong2023MetaGPT} are canonical examples,
explicitly modeling the waterfall SDLC: design $\rightarrow$
coding $\rightarrow$ testing. FlowGen~\cite{Lin2025Soen101}
operationalizes three SDLC models as distinct topologies:
FlowWater (strict waterfall chain), FlowTDD (requirement
$\rightarrow$ design $\rightarrow$ test $\rightarrow$
implementation $\rightarrow$ fix, a test-driven reordering), and
FlowScrum (cyclic iterative sprints). This paper is unique in
directly comparing the implications of different
SDLC-mirroring topologies for code quality.
L2MAC~\cite{Holt2023L2MAC} also follows a chain topology but with
a novel twist: each step in the instruction plan is executed by a
fresh-context agent, making the chain a sequence of independent
LLM invocations sharing only the external file store.

\textbf{\textit{Cyclic (Agile/Iterative) topologies}} introduce
back-edges that allow code to be revised in response to
verification feedback. AgentCoder~\cite{huang2023agentcoder}
implements a programmer $\rightarrow$ test executor $\rightarrow$
(if fail) $\rightarrow$ programmer cycle, bounded at 5
iterations. Self-Collaboration~\cite{Dong2024SelfCollaboration}
embeds a coder $\leftrightarrow$ tester back-edge within its
waterfall chain, max 4 iterations.
PairCoder~\cite{Zhang2024PairProgramming} enhances the cyclic
pattern with multi-plan exploration: a pool of $n$ solution plans
is pre-generated via k-means++ clustering for diversity, and the
cycle can switch to the next candidate plan when dead-end is
detected through history-based loop analysis.
MAGE~\cite{Zhao2024MAGE} combines a linear initialization chain
with a cyclic debug-judge loop, and introduces high-temperature
candidate sampling to explore multiple program variants
simultaneously.

\textbf{\textit{Hierarchical topologies}} place one or more
manager agents above a pool of worker agents, enabling
decomposition-and-delegation patterns. MAGIS~\cite{Tao2024Magis}
has a Manager that dynamically instantiates one Developer agent
per candidate file at runtime; each Developer edits its assigned
file and reports back to the manager-review layer. HyperAgent~\cite{Phan2024HyperAgent}
uses a planner above multiple repository navigation and editing
workers, combining top-down decomposition with bottom-up
repository evidence. SoA~\cite{Ishibashi2024SelfOrganized}
pushes this hierarchy further by allowing Mother agents to spawn
Child agents recursively according to inferred subtask
complexity. These systems treat harness orchestration itself as a
resource-allocation problem.

\textbf{\textit{Star topologies}} center on a hub agent that
coordinates multiple parallel worker agents. The
CANDOR~\cite{Xu2025Hallucination} Stage 3 panel is an example: a
Requirement Engineer fans out to three independent
Panelist+Interpreter pipelines, and the Curator aggregates their
outputs. MetaGPT~\cite{Hong2023MetaGPT}'s publish-subscribe
message pool creates a de facto star topology where the shared
pool serves as the hub.

\paragraph{Objective-driven and Adaptive Topologies}
A smaller but rapidly growing class of systems treats the
topology itself as a design variable to be optimized toward a
code quality signal. Recent systems such as
FlowReasoner~\cite{gao2025flowreasoner} and
BOAD~\cite{xu2025boad} further reinforce this trend by treating
multi-agent organization itself as an adaptive object to be
generated, searched, or optimized per task.

\textbf{\textit{Dynamic agent pool scaling}} is the simplest form
of adaptivity: the number of agents scales with task complexity,
but the topology type is fixed.
SoA~\cite{Ishibashi2024SelfOrganized} implements this via a
hierarchical tree of Mother and Child agents, where Mother agents
decide at runtime how many subfunctions to decompose into,
spawning corresponding Child agents. The key insight is that each
agent's context window remains bounded, as complexity is handled
by growing the agent pool rather than growing individual context
windows. MAGIS~\cite{Tao2024Magis} similarly instantiates
Developer agents dynamically based on the number of candidate
files identified during repository analysis. BOAD~\cite{xu2025boad} extends this line of thought from dynamic
scaling to hierarchy discovery: instead of manually fixing the
specialized sub-agent structure, it formulates the selection of
helpful localization, editing, and validation sub-agents as a
bandit-optimization problem, showing that automatically
discovered hierarchical teams can outperform manually designed ones.

\textbf{\textit{Feedback-driven DAG restructuring}} is best
represented by EvoMAC~\cite{Hu2025EvoMAC}. Its workflow is a DAG
whose nodes correspond to agents and whose edges define
information flow. After each iteration, a Gradient Agent reads
execution logs to attribute failures to agents, and an Updating
Agent modifies the prompts and graph structure. This is the only
system in the survey where the harness topology is structurally
modified in response to execution feedback.

\textbf{\textit{Runtime self-reorganization}} is
SEW~\cite{Liu2025SEW}'s approach: the system generates and mutates
entire workflow specifications using Direct Evolution (DE) and
Hyper Evolution (HE) operators applied to LLM-generated workflow
descriptions in structured formats (BPMN, CoRE, Python, YAML).
Rather than optimizing agent parameters, SEW~\cite{Liu2025SEW}
optimizes the workflow structure including the sequence of agent
calls, the routing logic, and the feedback paths. The two
canonical topologies it discovers (a linear chain and a feedback
loop) emerge from optimization rather than being hand-designed. FlowReasoner~\cite{gao2025flowreasoner} pushes this adaptive view
further by training a query-level meta-agent that generates a
tailored multi-agent system for each input problem under external
execution feedback, making topology selection itself part of the
deliberative inference process rather than a fixed system design.

\subsection{Execution Feedback and Shared-Harness Synchronization}

We discuss how a group of agents can exploit the executability of
code, and how they maintain a consistent shared view of the
program state.
This dimension is the defining one for code-centric MAS: the
shared harness is uniquely executable and produces objective
oracle signals. We address two sub-questions: what types of
execution feedback are used, and how is shared state
synchronized across agents.


\begin{table}[t]
\centering
\caption{Representative MAS execution-feedback and convergence designs.}
\label{tab:mas_feedback_systems}
\renewcommand{\arraystretch}{1.2}
\setlength{\tabcolsep}{5pt}
\footnotesize
\begin{tabularx}{\textwidth}{@{}
  >{\raggedright\arraybackslash}p{2.6cm}
  >{\raggedright\arraybackslash}p{3.8cm}
  >{\raggedright\arraybackslash}p{2.3cm}
  >{\raggedright\arraybackslash}p{3.3cm}
  >{\raggedright\arraybackslash}X@{}}
\toprule
\textbf{System} & \textbf{Harness Substrate} & \textbf{Topology} & \textbf{Execution Feedback} & \textbf{Convergence} \\
\midrule
\multicolumn{5}{@{}l}{\textit{Pre-defined topology}} \\
\addlinespace[2pt]
AgentCoder~\cite{huang2023agentcoder} & Execution & Cyclic & Test pass/fail & Correctness (test-gated) \\
MAGE~\cite{Zhao2024MAGE} & Execution (waveform) & Chain-cyclic & Checkpoint waveform & Score-based correctness \\
MapCoder~\cite{islam2024mapcoder} & Execution, implicit & Cyclic & Test pass/fail & Correctness \\
AutoSafeCoder~\cite{Nunez2024AutoSafeCoder} & Execution (static, fuzzer) & Cyclic & CWE warnings, crashes & Security convergence \\
QualityFlow~\cite{Hu2025QualityFlow} & Execution (real, imagined) & Gated cyclic & Pass/fail, imagined exec & Correctness (quality-gated) \\
CodeCoR~\cite{Pan2025CodeCoR} & Execution, implicit & Cyclic & Syntax, test pass/fail & Score-based soft correctness \\
MARCO~\cite{Rahman2025MACRO} & Execution (performance) & 2-node Cyclic & Time, memory, FLOPS & Performance, correctness \\
\midrule
\multicolumn{5}{@{}l}{\textit{Adaptive topology}} \\
\addlinespace[2pt]
SoA~\cite{Ishibashi2024SelfOrganized} & Execution, implicit gap & Hierarchical tree & Test pass/fail & Correctness (implicit fallback) \\
SEW~\cite{Liu2025SEW} & Implicit & Evolution & Test pass/fail & Implicit \\
EvoMAC~\cite{Hu2025EvoMAC} & Execution & Text DAG & Compiler, execution logs & Correctness (fixed-iteration) \\
FlowReasoner~\cite{gao2025flowreasoner} & Execution, implicit & Query workflow & Execution feedback & Objective-driven adaptive \\
Trae Agent~\cite{gao2025traeagent} & Repository, execution & Search pipeline & Test, pruning signals & Score-/selection-based \\
\bottomrule
\end{tabularx}
\end{table}

\subsubsection{Execution Feedback Integration}
\paragraph{Compiler and syntax feedback} Compiler and syntax
feedback catch structural errors before runtime and are used by
many systems. ChatDev~\cite{Qian2023ChatDev} feeds compiler
errors from the testing phase back to the programmer, though only
as one-off corrections within a single phase.
L2MAC~\cite{Holt2023L2MAC} runs syntax checks via its evaluator
module $E(D)$ after every file write, treating them as blocking
conditions that prevent the instruction pipeline from advancing.

\paragraph{Test pass/fail signals} Test pass/fail signals are the
most commonly used execution-feedback type.
AgentCoder~\cite{huang2023agentcoder} centers its entire loop on
whether independently generated test cases pass; the iteration
terminates on full pass or at the 5-iteration budget.
QualityFlow~\cite{Hu2025QualityFlow} introduces a notable
variant: Imagined Execution, in which an LLM simulates the Python
interpreter step-by-step and predicts test outcomes without
actually running the code, achieving 98\%+ precision and recall
on MBPP while avoiding label leakage from visible test cases.
The near-identical performance of
Self-Collaboration~\cite{Dong2024SelfCollaboration}'s simulated
LLM tester and its real-compiler ablation raises a provocative
empirical question: when is actual execution necessary, and when
can linguistic simulation of execution suffice?

\paragraph{Fuzzer crash traces} Fuzzer crash traces represent a
qualitatively different type of feedback: rather than a pass/fail
outcome, they provide a concrete failing input.
AutoSafeCoder~\cite{Nunez2024AutoSafeCoder} uses type-aware
mutation to generate crash-inducing input seeds and passes the
crashing input plus exit code to the Coding Agent. This
adversarial feedback is more informative than a generic failure
signal because it localizes the bug to a specific input category.

\paragraph{Static analysis warnings} Static analysis warnings
provide feedback about code structure and security properties
without execution. AutoSafeCoder~\cite{Nunez2024AutoSafeCoder}
uses CWE-mapped static analysis against the MITRE vulnerability
database, enabling the Static Analyzer Agent to suggest
remediation strategies keyed to specific vulnerability classes.

\paragraph{Performance profiling results} Performance profiling
results are uniquely exploited by MACRO~\cite{Rahman2025MACRO},
which treats code optimization as the primary task rather than
correctness. The Performance Evaluator Agent measures execution
time, memory usage, and FLOPS, and MACRO~\cite{Rahman2025MACRO} uniquely augments this
with real-time web search to retrieve relevant optimization
techniques from the research literature.

\paragraph{Fine-grained simulation feedback} MAGE~\cite{Zhao2024MAGE}'s
distinctive contribution is the finest-grained execution feedback
in the surveyed literature. Rather than reporting only whether a
testbench passes or fails, the State Checkpoint mechanism records
signal values at every clock edge and delivers to the Debug Agent
a waveform window around the first failing clock cycle. This
enables targeted repair at sub-test granularity.


\subsubsection{Shared-Harness Synchronization}

Sequential handoff is the most common synchronization mechanism:
each agent receives the output of its predecessor and passes its
own output to its successor. The program state exists only in the
form of the most recent artifact in the pipeline. This is
sufficient for simple linear pipelines but creates invisible
state divergence in multi-agent settings where multiple agents
modify the codebase in parallel or iteratively. {It
is also where the limits of code-mediated coordination become
clear. Even when agents share executable artifacts, the harness
still imposes information-theoretic constraints: channels have
finite bandwidth, summaries introduce compression loss, logs
become noisy, cached views go stale, and parallel branches raise
unresolved questions of authority and consistency. Code provides
a richer substrate for coordination, but it does not remove
these distributed-systems constraints.}

\paragraph{Shared blackboard} Shared blackboard provides a
globally accessible program state that all agents can read and
update. L2MAC~\cite{Holt2023L2MAC} implements this most cleanly:
the file store $D$ is an external, persistent structure that is
never overwritten but extended and revised. The Control Unit
manages all reads and writes, ensuring that each agent invocation
receives a precisely controlled context window.
MAGIS~\cite{Tao2024Magis}'s repository evolution memory $M$ is a
persistent key-value store mapping file versions to
LLM-generated summaries, updated incrementally via a specialized
blackboard for repository-level reasoning.
Self-Collaboration~\cite{Dong2024SelfCollaboration} is among the
first systems to explicitly name and invoke the blackboard
architecture, establishing a shared memory from which all three
roles read and write.

\paragraph{Parallel branches with merge} Parallel branches with
merge arise when multiple agents modify independent components
simultaneously, with their changes integrated at a later stage.
MAGIS~\cite{Tao2024Magis} instantiates one Developer per
candidate file; each modifies its assigned file independently,
and all changes are merged into the final repository diff.
HyperAgent~\cite{Phan2024HyperAgent} runs multiple Navigator and
Editor instances in parallel via Redis queues, with results
merged at the Planner level.

\paragraph{Structured context scheduling} Structured context
scheduling is the explicit management of what each agent sees and
when. It is the primary innovation of
L2MAC~\cite{Holt2023L2MAC}. The Control Unit resets the context
window between instruction steps, providing each new invocation
with a targeted summary of prior progress $(M_{rs})$ rather than
the full conversation history. When the context window approaches
capacity, the Control Unit stores partial results to $D$ and
re-initializes with a compressed view, explicitly instructing the
LLM which files to read or skip given the remaining context
margin. This mechanism solves the context-window problem not by
expanding the window but by carefully controlling its contents.
MetaGPT~\cite{Hong2023MetaGPT} implements a lighter form of
context scheduling via a publish-subscribe message pool: each
agent subscribes only to the document types relevant to its role,
receiving a filtered view of the shared state.

\paragraph{Hierarchical memory} Hierarchical memory combines
short-term working context with longer-term accumulated
knowledge. ChatDev~\cite{Qian2023ChatDev} explicitly separates
short-term memory (full dialogue within a phase) from long-term
memory (extracted solutions carried across phases).
Cogito~\cite{Li2025Cogito} implements hierarchical memory, drawing on
neurobiological architecture: short-term memory for immediate
task state, a long-term knowledge base for accumulated expertise,
and growth units for evolving abstractions that improve over
time. HyperAgent~\cite{Phan2024HyperAgent} uses a lightweight
LLaMA-3.1-8B summarizer to condense execution logs before storing
them in hierarchical memory, preventing context bloat.

\paragraph{Agent pool scaling} Agent pool scaling addresses the
context-management problem orthogonally: rather than managing
what a single agent sees, it distributes the context load across
more agents. SoA~\cite{Ishibashi2024SelfOrganized} is the
canonical example: by spawning more agents as task complexity
grows, each agent's context remains bounded. This is a
structural solution to the harness-state problem: instead of
building a shared representation that all agents can query,
SoA~\cite{Ishibashi2024SelfOrganized} partitions the task state
across agents, each holding a bounded slice. The limitation is
that global consistency is sacrificed: agents cannot reason about
the full program, only their assigned sub-function.

\paragraph{Other}
QualityFlow~\cite{Hu2025QualityFlow}'s revert mechanism
represents a synchronization pattern: the initial code artifact is never overwritten,
enabling the system to roll back to a prior shared harness state
if the debugging trajectory degrades quality. This is the only
work among the surveyed system that explicitly manages state history rather than always
moving forward.

\subsection{Position: The Shared Code-Centric Harness Substrate}

We propose a new position for the next generation of multi-agent
intelligence: the shared code-centric harness substrate. This
position is motivated by the central gap identified in the
literature: the lack of formal, persistent representations of the
shared code state that agents can query and update across
iterations. We argue that building such a harness substrate is
both feasible and necessary for achieving robust, scalable
multi-agent intelligence.



\begin{table}[t]
\centering
\caption{Representative MAS designs centered on shared program-state representation and synchronization.}
\label{tab:mas_world_state_systems}
\renewcommand{\arraystretch}{1.16}
\setlength{\tabcolsep}{4pt}
\footnotesize
\begin{tabularx}{\textwidth}{@{}
  >{\raggedright\arraybackslash}p{2.55cm}
  >{\raggedright\arraybackslash}p{3.05cm}
  >{\raggedright\arraybackslash}p{3.15cm}
  >{\raggedright\arraybackslash}p{3.05cm}
  >{\raggedright\arraybackslash}X@{}}
\toprule
\textbf{System}
& \textbf{Harness Substrate}
& \textbf{Agent Roles}
& \textbf{Execution Feedback}
& \textbf{Convergence / Synchronization} \\
\midrule
L2MAC~\cite{Holt2023L2MAC}
& Blackboard, repository, execution
& Plan, Synth, Verif (evaluator)
& Syntax, test pass/fail
& Correctness per instruction step \\ \midrule
Cogito~\cite{Li2025Cogito}
& Blackboard (3-tier memory)
& Neurobiological model
& NA
& Hierarchical memory synchronization \\ \midrule
CleanAgent~\cite{Qi2024CleanAgent}
& Execution (weak), implicit
& Plan, Understand, Synth, Exec
& Runtime errors
& Correctness through execution success \\ \midrule
Lingma SWE-GPT~\cite{Ma2024Lingma}
& Repository, execution
& Understand, Synth-Verif
& Syntax, git apply, tests
& Fixed-limit implicit convergence \\ \midrule
SyncMind~\cite{Guo2025SyncMind}
& Repository, execution (formal $S_k/B_k$)
& Synth-Understand, oracle Understand
& Test pass/fail, runtime errors
& Correctness, resource-constrained synchronization \\ \midrule
BOAD~\cite{xu2025boad}
& Repository, execution
& Orchestrator with specialized sub-agents
& Test pass/fail, validation reward
& Hierarchy discovery, coordination \\ \midrule
CANDOR~\cite{Xu2025Hallucination}
& Execution (Java, JaCoCo)
& Plan, Synth, Verif, Understand, Debate
& Compiler, coverage, tests
& Correctness, coverage, consensus \\
\bottomrule
\end{tabularx}
\end{table}


\subsubsection{Shared Harness Representation}

A foundational question for any MAS is: what is the substrate
these agents inhabit? In code as agent harness, the natural
answer is the shared program environment, namely the collection
of artifacts, execution contexts, and quality signals that agents
collectively act upon and that evolve as agents produce, revise,
and evaluate code. We call this the shared harness substrate, and
we distinguish four levels of formalization with which existing
systems represent it.

\paragraph{Implicit / File-only Representation}

The most common and least formalized category treats the shared
harness as simply the current code file or set of code files.
Agents receive the latest code artifact as part of their input
context and produce a modified or evaluated version. There is no
persistent, queryable representation: the shared state is
reconstructed implicitly at each agent invocation from the
conversational history. This category encompasses many
foundational systems: ChatDev~\cite{Qian2023ChatDev},
MetaGPT~\cite{Hong2023MetaGPT}, FlowGen~\cite{Lin2025Soen101},
MapCoder~\cite{islam2024mapcoder},
CodeCoR~\cite{Pan2025CodeCoR}, SEW~\cite{Liu2025SEW}, and
CodePori~\cite{Rasheed2024Codepori}. While this representation is
simple to implement, it entails a fundamental limitation: agents
cannot reason about the shared substrate except through the
narrow lens of their most recent context window.
State divergence~\cite{Guo2025SyncMind}, in which an agent's
internal belief about the code state diverges from the true
state, is invisible to the system and cannot be detected or
corrected.

\paragraph{Repository-based Representation}

A richer class of systems represents the shared harness as a
navigable repository: a file system with directory structure,
inter-file dependency graphs, call hierarchies, and version
history. This representation supports agents that reason about
where in the codebase a change needs to be made, what other
components depend on the changed function, and how the codebase
has evolved over time. MAGIS~\cite{Tao2024Magis} introduces a
repository evolution memory that caches file-level summaries and
incrementally updates them via git diff as files change across
issue-resolution episodes. HyperAgent~\cite{Phan2024HyperAgent}
provides agents with repository navigation tools
(get\_tree\_structure, go\_to\_definition, code\_search,
get\_all\_references), treating the repository as a structured
knowledge base. Lingma SWE-GPT~\cite{Ma2024Lingma} compresses the
repository view via abstract syntax tree (AST) skeletons,
preserving function signatures and class definitions to enable
efficient navigation. SyncMind~\cite{Guo2025SyncMind} is the only
work to formally define the repository substrate as a ground-truth
state $S_k$ and measure the divergence between $S_k$ and an
agent's belief state $B_k$.

\paragraph{Execution-based Representation}

Execution-based representation is the most distinctive category
for code generation. It has no direct parallel in general MAS and
represents the shared substrate through execution behavior. The
state is not what the code looks like but what the code does:
whether it compiles, which tests it passes, what vulnerabilities
a fuzzer uncovers, how fast it runs, and whether its runtime
behavior matches its specification. This execution-based
representation provides an objective oracle signal, a ground
truth that is not subject to the hallucination or bias that
affects purely linguistic agent evaluations. Systems that exploit
this representation include
AgentCoder~\cite{huang2023agentcoder},
AutoSafeCoder~\cite{Nunez2024AutoSafeCoder},
QualityFlow~\cite{Hu2025QualityFlow},
MACRO~\cite{Rahman2025MACRO}, EvoMAC~\cite{Hu2025EvoMAC},
CANDOR~\cite{Xu2025Hallucination}, and
MAGE~\cite{Zhao2024MAGE}. Notably,
MAGE~\cite{Zhao2024MAGE} achieves the finest-grained execution
feedback in the literature, operating at clock-edge granularity
via \textit{State Checkpoint} waveform snapshots.

\paragraph{Blackboard / Shared-State Representation}

A fourth category introduces an explicit, globally accessible
data structure that all agents can read from and write to (akin
to the classical blackboard architecture in
AI~\cite{erman1980hearsay}). This shared state is the closest
approximation in the literature to a formal harness substrate: it
persists across agent invocations, can be queried and updated,
and provides a consistent view of the program state to all
agents. Self-Collaboration~\cite{Dong2024SelfCollaboration} is
among the first systems to explicitly invoke the blackboard
metaphor, establishing a shared memory from which all three roles
(Analyst, Coder, Tester) read and write.
L2MAC~\cite{Holt2023L2MAC} implements the most principled
blackboard in the literature: a persistent file store $D$ with
semantically meaningful paths, accessed through a Control Unit
that explicitly manages which slice of state each agent
invocation sees. GameGPT~\cite{chen2023gamegpt} uses a shared
context buffer to reduce redundant information retransmission in
multi-round game development. Cogito~\cite{Li2025Cogito} draws on
neurobiological architecture to implement a three-tier memory:
short-term working state, long-term knowledge base, and growth
units for evolving abstractions, as a structured harness
representation.

\paragraph{The Central Gap}

The distribution of systems across these four categories reveals
a striking pattern: the majority of the literature resides in the
implicit/file-only category, lacking any formal model of the
shared harness substrate. This is the central gap that motivates
the code as agent harness framing. The program, uniquely among
multi-agent domains, is an artifact that executes. It produces
objective, non-linguistic signals that could in principle anchor
a formal shared substrate. Yet most systems fail to exploit this
property at the architectural level, instead relying on agents to
reason about code quality through natural language alone.


\subsubsection{Harness-State Convergence}


Convergence determines when a multi-agent coding harness should
stop iterating and accept its current program state as a
satisfactory outcome. In many existing MAS, convergence is still
defined implicitly, either by consensus among agents or by an
external iteration budget. However, code as agent harness has a
distinctive advantage: because the shared substrate is
executable, convergence can be grounded in objective behavioral
signals rather than in conversational agreement alone. We identify six convergence patterns, ranging from widely used test-gated and implicit convergence to less common security-, performance-, and consensus-based criteria.

\paragraph{Correctness convergence} Correctness convergence
(test-gated) is the most principled and widely used objective
criterion: the system terminates successfully when all test cases
pass. AgentCoder~\cite{huang2023agentcoder},
L2MAC~\cite{Holt2023L2MAC}, SyncMind~\cite{Guo2025SyncMind}, and
CANDOR~\cite{Xu2025Hallucination} implement test-gated
convergence. PairCoder~\cite{Zhang2024PairProgramming} augments
this with dead-end detection: if the same buggy code or feedback
appears in the iteration history, the system switches to the next
candidate plan rather than looping. FlowGen~\cite{Lin2025Soen101}
uses test-gated convergence but on LLM-generated tests rather
than ground-truth tests, introducing a potential quality concern:
a system can converge on code that passes its own biased tests
but fails on external evaluation.

\paragraph{Security convergence} Security convergence is uniquely
implemented by AutoSafeCoder~\cite{Nunez2024AutoSafeCoder}: the
system terminates successfully when no CWE vulnerabilities are
flagged by static analysis and no crashes are induced by the
fuzzer. This multi-criteria convergence is a strong argument for
the execution-based harness framing. Both convergence criteria
are grounded in objective program behavior, not agent opinions.

\paragraph{Performance convergence} Performance convergence is the
focus of MACRO~\cite{Rahman2025MACRO}: the optimization loop
terminates when user-defined runtime and memory thresholds are
satisfied, as measured by the Performance Evaluator against
actual execution benchmarks. This is the only system that treats
performance as the primary convergence criterion rather than
correctness.

\paragraph{Score-based convergence} Score-based convergence uses
quantitative quality scores computed by agents evaluating
intermediate outputs to determine when to stop.
MAGE~\cite{Zhao2024MAGE} ranks candidate programs by their
simulation mismatch score $s(r) = 1 - m(r) / tc(r)$ and
continues iterating until the maximum score reaches 1.0.
CodeCoR~\cite{Pan2025CodeCoR} uses a four-criteria binary score
(clarity, relevance, conciseness, context) to prune intermediate
outputs at each agent stage and selects the highest-ranked code
in its Ranked Code Set as the final output. It sets a soft
correctness convergence that submits the best available result
rather than waiting for a perfect solution. Trae Agent~\cite{gao2025traeagent} introduces a closely related
search-and-selection view at repository scale: it formulates
issue resolution as an optimal solution search problem and uses
modular generation, pruning, and selection agents to navigate a
large ensemble space of candidate patches. In this setting,
convergence is not only a matter of repeated repair, but also of
ranking, filtering, and selecting among competing solutions under
repository-aware evidence.

\paragraph{Consensus convergence} Consensus convergence aggregates
judgments from multiple reviewer agents.
CANDOR~\cite{Xu2025Hallucination} implements majority voting
among three Panelists on oracle correctness.
MAGIS~\cite{Tao2024Magis} uses LLM-judgment from the QA Engineer
as the acceptance signal, though this is a single-agent consensus
rather than a multi-agent vote. QualityFlow~\cite{Hu2025QualityFlow}
uses its Code Quality Checker as the single gating signal. It is
an efficient design where the quality checker serves as both a
convergence oracle and the system controller, enabling early exit
(75--84\% of problems converge after the first generator call).

\paragraph{Implicit convergence} Pipeline termination after a
fixed number of stages or iterations with no objective quality
criterion is the most prevalent convergence pattern in the
literature and represents the most significant gap in the field.
ChatDev~\cite{Qian2023ChatDev} terminates after a fixed number of
phases, or when two consecutive rounds produce identical code, or
after 10 rounds, none of which is an objective quality signal.
MetaGPT~\cite{Hong2023MetaGPT} terminates after completing the
fixed SOP stages.
Self-Collaboration~\cite{Dong2024SelfCollaboration} falls back to
implicit convergence after $n = 4$ iterations if the tester never
approves. EvoMAC~\cite{Hu2025EvoMAC} runs a fixed $K$ iterations
of the textual backpropagation loop. The prevalence of implicit
convergence is a direct consequence of the lack of formal shared
substrates: without an objective representation of the program
state, systems have no principled criterion for convergence.


\subsection{Patterns and Trends}

Across systems, differences in role specialization, shared-state representation,
execution grounding, and workflow topology are not independent
engineering choices; they interact to determine how reliably a
group of agents can maintain coherence over long-horizon coding
tasks. In this subsection, we distill the main
trends that emerge from the surveyed systems, highlighting both the common
structural bottlenecks of current systems and the design
principles that point toward more robust shared harnesses.


\paragraph{The implicit-harness-state constraint} The majority of
surveyed systems (ChatDev~\cite{Qian2023ChatDev},
MetaGPT~\cite{Hong2023MetaGPT}, FlowGen~\cite{Lin2025Soen101},
CodePori~\cite{Rasheed2024Codepori}, SEW~\cite{Liu2025SEW},
MapCoder~\cite{islam2024mapcoder},
CodeCoR~\cite{Pan2025CodeCoR}) operate without explicit
representations of the shared code harness. These systems rely on
agents to reconstruct state implicitly from conversational
history at each invocation. This design choice works for
function-level tasks where the program state is simple and does
not fragment across agents. However, this implicit approach
creates a fundamental vulnerability: without a formal shared
substrate, agents cannot reliably detect when their internal
understanding diverges from the true program
state~\cite{Guo2025SyncMind}. From the code as agent harness
perspective, the reliance on implicit state representations is
the technical root of system brittleness rather than a
scalability convenience.

\paragraph{Code-mediated channels do not eliminate coordination bottlenecks}
{The shift from free-form dialogue to code-mediated
coordination is a genuine architectural advance, but it should
not be overstated. Files, APIs, diffs, tests, logs, schemas,
blackboards, and workflow states are all partial channels
through which task state is encoded, transmitted, and
reconstructed. Each channel trades off fidelity, latency, and
scope: tests compress semantics into pass/fail, summaries save
context at the cost of detail, logs are grounded but noisy, and
shared blackboards improve persistence while creating authority
and consistency problems. The central design question is
therefore not merely whether code is present, but which
artifacts are authoritative, how they are compressed, and how
conflicts across channels are resolved.}

\paragraph{Execution feedback as the bridge between linguistic and formal reasoning}
The deepest divide in the literature is between systems that use
execution as ground truth and those that rely on linguistic model
judgments. Systems that ground shared state in execution
(AgentCoder~\cite{huang2023agentcoder},
AutoSafeCoder~\cite{Nunez2024AutoSafeCoder},
QualityFlow~\cite{Hu2025QualityFlow}, EvoMAC~\cite{Hu2025EvoMAC},
MAGE~\cite{Zhao2024MAGE}) have access to objective oracle
signals, signals that cannot hallucinate. Yet a surprising
finding complicates this picture:
Self-Collaboration~\cite{Dong2024SelfCollaboration} and
QualityFlow~\cite{Hu2025QualityFlow} demonstrate that
LLM-simulated execution can achieve 98\%+ precision and recall in
predicting actual outcomes without running code. This suggests
that execution feedback's value is not uniform across all failure
modes. It excels at detecting the corner cases that linguistic
simulation structurally cannot imagine (runtime crashes, resource
exhaustion, boundary condition errors, performance regressions),
but for many correctable bugs, simulated reasoning may suffice. A
mature harness would integrate both: using linguistic reasoning
as the fast path and delegating to execution as the verification
oracle only for the failure modes that require it.

\paragraph{Two complementary representations of the shared harness}
The surveyed systems reveals two conceptually orthogonal views:
repository-based representation (structure: what functions call
what, where does data flow, what are the dependencies) and
execution-based representation (behavior: what does the code do
when run, how does state evolve at runtime, what emergent
failures occur under different inputs). MAGIS~\cite{Tao2024Magis}
and HyperAgent~\cite{Phan2024HyperAgent} operate primarily in the
repository view, enabling agents to reason about codebase
architecture. AgentCoder~\cite{huang2023agentcoder} and
MAGE~\cite{Zhao2024MAGE} operate primarily in the execution view,
grounding shared state in runtime signals. Yet none of the
surveyed systems fully unifies both views into a single harness
substrate where agents can reason across both the static
structure of code and its dynamic behavior. The deepest harness
would integrate these two perspectives, answering questions like
``which components are slow'' (requires both call graphs and
profiling data) or ``does this refactoring break APIs that
external code depends on'' (requires both static analysis and
dynamic testing).

\paragraph{Topology complexity inversely correlates with harness-state formality}
Systems with explicit, formal shared substrates use simpler
topologies, while systems lacking formal shared state employ
increasingly complex topology patterns as a structural
workaround. L2MAC~\cite{Holt2023L2MAC}, which has the clearest
formal harness substrate (a persistent file store with explicit
context scheduling), uses a simple sequential chain with
sophisticated state management. By contrast,
implicit-state systems like EvoMAC~\cite{Hu2025EvoMAC} and
SEW~\cite{Liu2025SEW} develop elaborate adaptive topologies
(dynamic DAGs, workflow mutation, agent pool scaling) that
attempt to optimize the collaboration structure in the absence of
a principled shared representation. This suggests that topology
complexity is partially a symptom: when the substrate is
formally represented and queryable, agents can coordinate through
simple, transparent protocols. When the substrate is implicit,
agents require richer interaction patterns to compensate for
missing state information.

\paragraph{Context management is the tax of implicit shared state}
A striking pattern is that many systems have
developed sophisticated context-management mechanisms precisely
because they lack a formal shared substrate.
L2MAC~\cite{Holt2023L2MAC}'s Control Unit,
MetaGPT~\cite{Hong2023MetaGPT}'s publish-subscribe pool,
SoA~\cite{Ishibashi2024SelfOrganized}'s agent-pool scaling, and
Cogito~\cite{Li2025Cogito}'s three-tier memory are all responses
to the same underlying problem: how to give agents a coherent
view of a code harness that is too large to fit in any one
context window. A mature harness substrate could unify these
disparate solutions by providing a principled, queryable
representation of task state that agents access on demand,
rather than forcing the system to carefully manage what each
agent sees at every step.

\paragraph{Agent specialization increases the criticality of shared state metrics}
As agent role diversity increases, from basic coder-tester pairs
to systems with Architect, Manager, Navigator, Executor, and
Verifier roles, the need for a unified shared substrate becomes
urgent. Without shared understanding of code state, the Planning
Agent may decompose tasks based on an outdated codebase snapshot,
the Execution Agent may run tests against a different version
than the Synthesis Agent intended, and the Verification Agent's
feedback may misfire. EvoMAC~\cite{Hu2025EvoMAC} addresses this
through its Gradient and Updating agents that explicitly monitor
failure attribution at the MAS level.
SyncMind~\cite{Guo2025SyncMind} formalizes the problem as agent
belief divergence $|B_k - S_k|$, proposing explicit
synchronization protocols. The proliferation of agent roles is
thus not merely an engineering choice. It is a forcing function
for developing more mature shared harnesses. Multi-agent systems
with rich role repertoires cannot function robustly without them.
